% Chapter 4, Topic _Linear Algebra_ Jim Hefferon
%  http://joshua.smcvt.edu/linearalgebra
%  2020-Apr-08
\topic{计算机图形学}
\index{computer graphics|(}
前面关于射影几何的专题给出了这个模型，
描述我们的眼睛或照相机如何观看世界。
\begin{center} 
  \vcenteredhbox{\includegraphics{det/mp/ch4.14}}
\end{center}
过原点的一条直线上的所有点都投影到同一个
位置。

在那个专题中我们定义了：对任一非零向量 $\vec{v}\in\Re^3$，与之相伴的
\definend{射影平面中的点 $p$}\index{point!in projective plane} 
是集合 $\set{k\vec{v}\suchthat \text{$k\in\Re$ and $k\neq 0$}}$。
这是与 $\vec{v}$ 位于同一条过原点直线上的
非零向量的全体。

要描述一个射影点，我们可以给出该直线的任一代表元。
因此下面每一个都表示同一个射影点。
\begin{equation*}
  \colvec[r]{1 \\ 2 \\ 3}
  \qquad
  \colvec{1/3 \\ 2/3 \\ 1}
  \qquad
  \colvec[r]{-2 \\ -4 \\ -6}
\end{equation*} 
每一个都是点~$p$ 的一个
\definend{齐次坐标向量}\index{coordinates!homogeneous}%
\index{homogeneous coordinate vector}\index{vector!homogeneous coordinate}。 
两个齐次坐标向量（按定义非零）
\begin{equation*}
  \tilde{p}_1=\colvec{a_1 \\ b_1 \\ c_1}
  \qquad
  \tilde{p}_2=\colvec{a_2 \\ b_2 \\ c_2}
\end{equation*}
表示同一个射影点，当且仅当存在缩放因子
$s\neq 0$ 使得 $s\tilde{p}_1=\tilde{p}_2$。

在无穷多个可能的代表元中，
我们常使用第三个分量为~$1$ 的那一个。
这相当于投影到平面 $z=1$ 上。
\begin{center} 
  \vcenteredhbox{\includegraphics{det/mp/ch4.18}}
\end{center}

在本专题中，我们将展示如何利用这些想法来实现计算机图形学中的
一些效果。
为此，我们把前面的图重画为不带球，
在原点处放一台电影放映机，而平面 $z=1$ 看起来像
电影院的银幕。
\begin{center} 
  \vcenteredhbox{\includegraphics{det/mp/ch4.60}}
\end{center}
这将三维空间中灰线上的向量与银幕平面中的~$p$
联系起来。
\begin{equation*}
  \tilde{p}=\colvec{a \\ b \\ c} 
  \mapsto 
  p=\colvec{x \\ y}=\colvec{a/c \\ b/c}
\end{equation*}

我们可以改造已经见过的关于矩阵的内容来实现这些
变换。
旋转\index{rotation} 
就是一例。
这个矩阵使平面~$z=1$ 绕原点旋转角度~$\theta$。
\begin{equation*}
  \begin{mat}
    \cos\theta  &-\sin\theta  &0  \\
    \sin\theta  &\cos\theta   &0  \\
    0           &0            &1  
  \end{mat}
  \colvec{x \\ y \\ 1}
  =
  \colvec{cos\theta\cdot x -\sin\theta\cdot y  \\  
         \sin\theta\cdot x +\cos\theta\cdot y  \\ 
         1}
\end{equation*}
注意它对任何齐次坐标向量都适用；
如果我们先用这个矩阵
\begin{equation*}
  \begin{mat}
    \cos\theta  &-\sin\theta  &0  \\
    \sin\theta  &\cos\theta   &0  \\
    0           &0            &1  
  \end{mat}
  \colvec{a \\ b \\ c}
  =
  \colvec{cos\theta\cdot a -\sin\theta\cdot b  \\  
         \sin\theta\cdot a +\cos\theta\cdot b  \\ 
         c}
\end{equation*}
然后移到 $z=1$~平面
\begin{equation*}
  \colvec{cos\theta\cdot a -\sin\theta\cdot b  \\  
         \sin\theta\cdot a +\cos\theta\cdot b  \\ 
         c}
  \mapsto
  \colvec{(cos\theta\cdot a -\sin\theta\cdot b)/c  \\  
         (\sin\theta\cdot a +\cos\theta\cdot b)/c  \\ 
         1}
\end{equation*}
那么得到的结果与
先移到该平面再用
矩阵作用的结果相同。
\begin{equation*}
  \begin{mat}
    \cos\theta  &-\sin\theta  &0  \\
    \sin\theta  &\cos\theta   &0  \\
    0           &0            &1  
  \end{mat}
  \colvec{a/c \\ b/c \\ 1}
  =
  \colvec{cos\theta\cdot a/c -\sin\theta\cdot b/c  \\  
         \sin\theta\cdot a/c +\cos\theta\cdot b/c  \\ 
         1}
\end{equation*}

所以使用齐次坐标并无妨害。
但它的优势是什么呢？

计算机图形学中的平移操作，\index{translation} 
即把物体从一处滑动到另一处，
不是线性变换，因为它
不保持原点不动。
但如果使用齐次坐标，我们就可以用矩阵来实现。
这个矩阵把所讨论平面中的点沿 $x$~方向平移 $t_x$，沿 $y$~方向平移 $t_y$。
\begin{equation*}
  \begin{mat}
    1  &0  &t_x  \\
    0  &1  &t_y  \\
    0  &0  &1  
  \end{mat}
  \colvec{a \\ b \\ c}
  =
  \colvec{a+t_x\cdot c  \\  
          b+t_y\cdot c  \\ 
          c}
  \mapsto
  \colvec{a/c+t_x  \\  
          b/c+t_y  \\ 
          1}
\end{equation*}
也就是说，在所讨论的平面中，这个矩阵把
$\binom{a/c}{b/c}$ 滑动到 $\binom{a/c+t_x}{b/c+t_y}$。
这样，齐次坐标使我们能够使用矩阵。

好，那么使用这些矩阵的优势是什么？
多出来的这个坐标能带给我们什么？
假设我们正在用计算机制作一部电影。
此刻
摄影机正在同时平移和旋转。
场景中的每一个点都需要既被平移又被旋转。
与其让计算机对每个点执行两次操作，
我们可以把两个矩阵相乘，
这样计算机对每个点只需施加一次操作：把该点
乘以所得的矩阵。
这是巨大的加速与简化。

我们将列举一些可以实现的效果的例子。
前面已经谈过旋转。
这是旋转半个弧度的图示。
\begin{center} 
  \vcenteredhbox{\includegraphics{det/mp/ch4.61}}
  \quad$\mapsto$\quad
  \vcenteredhbox{\includegraphics{det/mp/ch4.62}}
\end{center}
这是 $t_x=1.5$、~$t_y=0.5$ 的平移。
\begin{center} 
  \vcenteredhbox{\includegraphics{det/mp/ch4.61}}
  \quad$\mapsto$\quad
  \vcenteredhbox{\includegraphics{det/mp/ch4.67}}
\end{center}

接下来是缩放。
这个矩阵把目标平面中的物体
沿 $x$~方向缩放~$s$ 倍，
沿 $y$~方向缩放~$t$ 倍。
\begin{equation*}
  \begin{mat}
    s  &0  &0  \\
    0  &t  &0  \\
    0  &0  &1  
  \end{mat}
  \colvec{a/c \\ b/c \\ 1}
  =
  \colvec{s\cdot a/c  \\  
          t\cdot b/c \\ 
          1}
\end{equation*}
在这幅图中，我们沿 $x$~方向缩放 $s=2.5$ 倍，
沿 $y$~方向缩放~$t=0.75$ 倍。
\begin{center} 
  \vcenteredhbox{\includegraphics{det/mp/ch4.61}}
  \quad$\mapsto$\quad
  \vcenteredhbox{\includegraphics{det/mp/ch4.63}}
\end{center}

如果取 $s=t$，那么整个形状被等比缩放。
例如，如果我们把 $s=t=1.10$ 的帧串联起来，那么在电影中
看起来物体正在向我们靠近。
\begin{center} 
  \vcenteredhbox{\includegraphics{det/mp/ch4.61}}
  \quad
  \vcenteredhbox{\includegraphics{det/mp/ch4.68}}
  \quad
  \vcenteredhbox{\includegraphics{det/mp/ch4.69}}
  \quad
  \vcenteredhbox{\includegraphics{det/mp/ch4.70}}
\end{center}

我们可以反射物体。
这是关于直线 $y=x$ 的反射。
\begin{equation*}
  \begin{mat}
    0  &1  &0  \\
    1  &0  &0  \\
    0  &0  &1  
  \end{mat}
  \colvec{a/c \\ b/c \\ 1}
  =
  \colvec{b/c  \\  
          a/c \\ 
          1}
\end{equation*}
这里的虚线是 $y=x$。
\begin{center} 
  \vcenteredhbox{\includegraphics{det/mp/ch4.61}}
  \quad$\mapsto$\quad
  \vcenteredhbox{\includegraphics{det/mp/ch4.64}}
\end{center}

这是关于 $y=-x$ 的反射。
\begin{equation*}
  \begin{mat}
    0   &-1  &0  \\
    -1  &0  &0  \\
    0   &0  &1  
  \end{mat}
  \colvec{a/c \\ b/c \\ 1}
  =
  \colvec{-b/c  \\  
          -a/c \\ 
          1}
\end{equation*}
下面的虚线是 $y=-x$。
\begin{center} 
  \vcenteredhbox{\includegraphics{det/mp/ch4.61}}
  \quad$\mapsto$\quad
  \vcenteredhbox{\includegraphics{det/mp/ch4.65}}
\end{center}

更复杂的变换也是可能的。
这是一个\definend{错切}。\index{shear}
\begin{equation*}
  \begin{mat}
    1   &1  &0  \\
    0   &1  &0  \\
    0   &0  &1  
  \end{mat}
  \colvec{a/c \\ b/c \\ 1}
  =
  \colvec{a/c+b/c  \\  
          b/c \\ 
          1}
\end{equation*}
在这幅图中，点的 $y$~分量不变，
但 $x$~分量加上了~$y$ 的值。
\begin{center} 
  \vcenteredhbox{\includegraphics{det/mp/ch4.61}}
  \quad$\mapsto$\quad
  \vcenteredhbox{\includegraphics{det/mp/ch4.66}}
\end{center}

这一切都用矩阵来表示的一大优势是，我们可以通过组合简单的操作来完成
复杂的操作。
要关于直线 $y=-x+2$ 作反射，我们可以找到三个矩阵：
先把所有东西滑动到原点，再关于 $y=-x$ 反射，然后
滑回去。
\begin{equation*}
  \begin{mat}
    1   &0  &0  \\
    0   &1  &2  \\
    0   &0  &1  
  \end{mat}
  \begin{mat}
    0   &-1  &0  \\
    -1  &0   &0  \\
    0   &0   &1  
  \end{mat}
  \begin{mat}
    1   &0  &0  \\
    0   &1  &-2  \\
    0   &0  &1  
  \end{mat}
  \tag{*}
\end{equation*}
（一如既往，先执行的操作由
右边的矩阵描述。
也就是说，右边的矩阵描述把所讨论平面中的所有点滑动 $-2$，
中间的矩阵关于 $y=-x$ 作反射，
而左边的矩阵把所有点滑回去。）

还有用矩阵可以实现的更复杂的效果。
下面是
\definend{一般仿射变换}\index{affine transformation} 
和
\definend{一般射影变换}\index{projective transformation}
的矩阵。
\begin{equation*}
  \begin{mat}
    d   &e  &f  \\
    g   &h  &i  \\
    0   &0  &1  
  \end{mat}
  \qquad
  \begin{mat}
    d   &e  &f  \\
    g   &h  &i  \\
    j   &k  &1  
  \end{mat}
\end{equation*}
不过，对它们几何效果的描述超出了本书的范围。

关于计算机图形学有大量文献，
线性代数在其中扮演着重要角色。
一份极好的参考资料是 \cite{Hughes}。
这门学科是数学与艺术的绝妙融合；
参见 \cite{Disney}。



% ========================================
\begin{exercises}
  \item 
    计算 ($*$) 中的乘积。
    \begin{answer}
      直接计算得乘积如下。
      \begin{equation*}
        \begin{mat}
          1   &0  &0  \\
          0   &1  &2  \\
          0   &0  &1  
        \end{mat}
        \begin{mat}
          0   &-1  &0  \\
          -1  &0   &0  \\
          0   &0   &1  
        \end{mat}
        \begin{mat}
          1   &0  &0  \\
          0   &1  &-2  \\
          0   &0  &1  
        \end{mat}
        =
        \begin{mat}
         0 &-1 &2 \\
        -1 &0  &2 \\
         0 &0  &1
        \end{mat}
      \end{equation*}
    \end{answer}
% sage: M0 = matrix(QQ, [[1,0,0], [0,1,2], [0,0,1]])
% sage: M0
% [1 0 0]
% [0 1 2]
% [0 0 1]
% sage: M1 = matrix(QQ, [[0,-1,0], [-1,0,0], [0,0,1]])
% sage: M2 = matrix(QQ, [[0,-1,0], [-1,0,0], [0,0,1]])
% sage: M2 = matrix(QQ, [[1,0,0], [0,1,-2], [0,0,1]])
% sage: M0*M1*M2
% [ 0 -1  2]
% [-1  0  2]
% [ 0  0  1]
  \item 求关于直线 $y=2x$ 作反射的矩阵。
    \begin{answer}
      在 $\R^2$ 中工作，设该矩阵为 $M$。
      我们得到这两个。
      \begin{equation*}
        \colvec{1 \\ 2}\mapsto \colvec{1 \\ 2}
        \qquad
        \colvec{-2 \\ 1}\mapsto \colvec{2 \\ -1}
      \end{equation*}
      （对于第二个，起始向量位于过原点且
      垂直于 $y=2x$ 的直线上。）
      我们有
      \begin{equation*}
        M
        \begin{mat}
          1  &-2  \\
          2  &1
        \end{mat}
        =
        \begin{mat}
          1  &2  \\
          2  &-1
        \end{mat}
      \end{equation*}
      求解得
      \begin{equation*}
        M =
            \begin{mat}
              1  &2  \\
              2  &-1
            \end{mat}
            \begin{mat}
              1  &-2  \\
              2  &1     
            \end{mat}^{-1}
          =
            \begin{mat}
              -3/5  &4/5 \\
              4/5   &3/5
            \end{mat}
      \end{equation*}
      于是对齐次坐标而言，矩阵是这个。
      \begin{equation*}
      \begin{mat}
        -3/5  &4/5  &0  \\
         4/5  &3/5  &0  \\
           0  &0    &1
      \end{mat}
      \end{equation*}
    \end{answer}
  \item 求关于直线 $y=2x-4$ 作反射的矩阵。
    \begin{answer}
      先把所有点移动 $4$，用前一题的结果关于直线 $y=2x$ 作反射，然后把它们移回去。
      \begin{equation*}
      \begin{mat}
        1  &0  &0  \\
        0  &1  &-4 \\
        0  &0  &1
      \end{mat}
      \begin{mat}
        -3/5  &4/5  &0  \\
         4/5  &3/5  &0  \\
           0  &0    &1
      \end{mat}
      \begin{mat}
        1  &0  &0  \\
        0  &1  &4 \\
        0  &0  &1
      \end{mat}
      =
      \begin{mat}
        -3/5  &4/5  &16/5  \\
         4/5  &3/5  &-8/5 \\
           0  &0    &1
      \end{mat}
      \end{equation*}
    \end{answer}
% sage: M0 = matrix(QQ, [[1,0,0], [0,1,-4], [0,0,1]])
% sage: M1 = matrix(QQ, [[-3/5,4/5,0], [4/5,3/5,0], [0,0,1]])
% sage: M2 = matrix(QQ, [[1,0,0], [0,1,4], [0,0,1]])
% sage: M0*M1*M2
% [-3/5  4/5 16/5]
% [ 4/5  3/5 -8/5]
% [   0    0    1]
  \item 
    旋转与平移都是刚体操作。
    先旋转再平移的矩阵是什么？
    \begin{answer}
    \begin{equation*}
      \begin{mat}
        1  &0  &t_x  \\
        0  &1  &t_y  \\
        0  &0  &1
      \end{mat}
      \begin{mat}
        \cos\theta &-\sin\theta  &0 \\
        \sin\theta &\cos\theta   &0 \\
        0          &0            &1
      \end{mat}
      =
      \begin{mat}
        \cos\theta  &-\sin\theta  &t_x  \\
        \sin\theta  &\cos\theta   &t_y  \\
        0           &0            &1
      \end{mat}
    \end{equation*}
    \end{answer}
  \item 齐次坐标以显然的方式推广到对三维空间的操作：~每一坐标
    是一组四个分量的非零向量，它们互为倍数而彼此关联。
    给出绕 $z$~轴旋转的矩阵，以及
    绕 $y$~轴旋转的矩阵。
    \begin{answer}
    下面是两个 $\nbyn{4}$ 矩阵。
    \begin{equation*}
      \begin{mat}
        \cos\theta  &-\sin\theta  &0  &0  \\
        \sin\theta  &\cos\theta   &0  &0  \\
        0           &0            &0  &0  \\
        0           &0            &0  &1  
      \end{mat}
      \qquad
      \begin{mat}
        \cos\theta  &0  &-\sin\theta  &0   \\
        0           &0  &0            &0  \\
        \sin\theta  &0  &\cos\theta   &0   \\
        0           &0  &0            &1  
      \end{mat}
    \end{equation*}    
    \end{answer}
\end{exercises}
\index{computer graphics|)}
\endinput
